% =========================================
% PHẦN 5: Kết quả thực nghiệm
% =========================================

\begin{frame}
  \centering
  \Huge \textbf{PHẦN 5} \\[0.5cm]
  \LARGE Kết quả thực nghiệm
\end{frame}

%---------------------------------
\begin{frame}{Thiết lập thực nghiệm}
\begin{columns}
\column{0.50\textwidth}
\textbf{Môi trường mô phỏng:}
\begin{itemize}
    \item Kịch bản: Ngã tư 4 hướng, Hà Nội mẫu (SUMO)
    \item Thời gian mỗi episode: 3600 giây (1 giờ)
    \item Warmup: 60 giây trước khi agent can thiệp
    \item Lưu lượng EW:NS = 2:1; 60\% xe máy
    \item GUI: tắt (headless) khi huấn luyện
\end{itemize}

\column{0.46\textwidth}
\textbf{Baseline so sánh — Fixed-Time 160s:}
\begin{itemize}
    \item EW xanh: 100s $\to$ EW vàng: 5s
    \item NS xanh: 50s $\to$ NS vàng: 5s
    \item Chu kỳ cố định, \textbf{không thích ứng}
    \item Tỉ lệ 2:1 đã được phản ánh trong baseline
\end{itemize}
\end{columns}

\vspace{8pt}
\textbf{5 chỉ số đánh giá chính:}
\begin{itemize}
    \item Độ dài hàng đợi trung bình và cực đại (xe dừng)
    \item Thời gian chờ trung bình (giây/bước)
    \item Tốc độ lưu thông trung bình (m/s)
    \item Lưu lượng xe qua nút / episode (throughput)
    \item Tổng phần thưởng (Total Reward)
\end{itemize}
\end{frame}

%---------------------------------
\begin{frame}{Kết quả so sánh DQN vs. Fixed-Time}
\begin{center}
\begin{tabular}{lcccr}
\toprule
\textbf{Chỉ số} & \textbf{DQN} & \textbf{Fixed-Time} & \textbf{Thay đổi} & \textbf{Cải thiện} \\
\midrule
Hàng đợi TB (xe)    & \textbf{16.05} & 33.13 & $-17.08$ & $\downarrow$51.6\% \\
Hàng đợi Max (xe)   & \textbf{51}    & 66    & $-15$    & $\downarrow$22.7\% \\
Độ lệch chuẩn       & \textbf{8.16}  & 12.77 & $-4.61$  & $\downarrow$36.1\% \\
Thời gian chờ (s)   & \textbf{30.55} & 121.35& $-90.80$ & $\downarrow$74.8\% \\
Tốc độ TB (m/s)     & \textbf{5.65}  & 4.87  & $+0.78$  & $\uparrow$16.0\%  \\
Throughput (xe)     & \textbf{344}   & 338   & $+6$     & $\uparrow$1.8\%   \\
Total Reward        & $\mathbf{-14.754}$ & $-35.512$ & $+20.758$ & $\uparrow$58.5\% \\
\bottomrule
\end{tabular}
\end{center}
\begin{block}{}
\centering DQN vượt trội hơn Fixed-Time trên tất cả 7 chỉ số
\end{block}
\end{frame}

%---------------------------------
\begin{frame}{Phân tích kết quả (1): Hàng đợi và Thời gian chờ}
\begin{itemize}
    \item \textbf{Hàng đợi trung bình giảm 51.6\%}
    \begin{itemize}
        \item DQN: 16.05 xe/làn \quad Fixed-Time: 33.13 xe/làn
        \item Độ lệch chuẩn thấp hơn 36.1\% → DQN \textbf{ổn định hơn}, ít đợt ùn tắc đột ngột
    \end{itemize}
    \item \textbf{Thời gian chờ giảm 74.8\%} — cải thiện lớn nhất trong tất cả chỉ số
    \begin{itemize}
        \item DQN: 30.55 s/bước \quad Fixed-Time: 121.35 s/bước
    \end{itemize}
    \item \textbf{Lý do:} Agent học phân bổ pha xanh đúng lúc, đúng hướng:
    \begin{itemize}
        \item Quan sát $q_i$, $w_i$ liên tục theo thời gian thực
        \item Tập trung thời gian xanh vào hướng đang tích lũy xe
        \item Không lãng phí thời gian xanh khi hướng đó vắng xe
    \end{itemize}
    \item \textbf{Fixed-Time không thể làm điều này:} dù NS tắc hay EW trống, đèn vẫn giữ đúng 100s/50s cố định.
\end{itemize}
\end{frame}

%---------------------------------
\begin{frame}{Phân tích kết quả (2): Tốc độ, Throughput và Reward}
\begin{itemize}
    \item \textbf{Tốc độ lưu thông tăng 16.0\%}
    \begin{itemize}
        \item DQN: 5.65 m/s \quad Fixed-Time: 4.87 m/s
        \item DQN không chỉ giảm hàng đợi mà còn giúp xe di chuyển nhanh hơn tổng thể
    \end{itemize}
    \item \textbf{Throughput tăng 1.8\%} (344 vs 338 xe/episode)
    \begin{itemize}
        \item Mức cải thiện khiêm tốn — cả hai chiến lược xử lý lượng xe tương đương
        \item Nhưng DQN thực hiện điều này với \textbf{ít ùn tắc và thời gian chờ thấp hơn đáng kể}
    \end{itemize}
    \item \textbf{Total Reward cải thiện 58.5\%}
    \begin{itemize}
        \item DQN: $-14.754$ \quad Fixed-Time: $-35.512$
        \item Reward âm tỉ lệ với mức ùn tắc → DQN cao hơn nghĩa là ít ùn tắc hơn
        \item Xác nhận agent đã tối ưu hóa thành công hàm mục tiêu đề ra
    \end{itemize}
    \item \textbf{Chiến lược học được:} giữ EW $\geq$ 60s, không bỏ quên NS, tránh penalty $-1.0$
\end{itemize}
\end{frame}

%---------------------------------
\begin{frame}{Kết luận và Hướng phát triển}
\textbf{Kết quả đạt được:}
\begin{itemize}
    \item Mô hình hóa thành công bài toán điều khiển đèn dưới dạng MDP với đặc thù giao thông Việt Nam
    \item Double DQN + Dueling Network vượt trội hơn Fixed-Time 160s trên tất cả các chỉ số
    \item Giảm hàng đợi \textbf{51.6\%}, thời gian chờ \textbf{74.8\%}, tăng tốc độ \textbf{16.0\%}
\end{itemize}

\vspace{6pt}
\textbf{Hướng phát triển:}
\begin{itemize}
    \item \textbf{Multi-Agent RL (MARL)}: điều phối đồng thời nhiều nút giao thông
    \item \textbf{Transfer Learning}: tái sử dụng model đã huấn luyện cho kịch bản mới
    \item \textbf{Dữ liệu thực}: tích hợp camera/cảm biến tại nút giao Hà Nội thực tế
    \item \textbf{Reward nâng cao}: bổ sung khí thải, nhiên liệu, mức độ an toàn
\end{itemize}
\end{frame}

%---------------------------------
\begin{frame}{}
\centering
\vspace{2cm}
{\LARGE \textbf{Cảm ơn thầy/cô và các bạn đã lắng nghe!}}\\[1cm]
\end{frame}
